\section{Successor Residue}

The extensional ledger says which finite obligations count as the same. It
does not yet say what happens when two gauge transports are both admissible
but do not commute. Gauge order matters because a transport is not a decorative
arrow between completed points. It is a rule for carrying a distinction from
one finite reading to the next. If the order of two transports changes the
carried distinction, the difference is a residue the grammar must read.

The usual notation for the commutator is short:
\[
  [A,B] = AB - BA .
\]
The notation is useful only after the grammar has earned it. The expression
\(AB\) says that the \(A\)-transport is carried and then the \(B\)-transport
is carried. The expression \(BA\) says the reverse. The bracket records the
debt left when those two ordered readings fail to identify. The residue is
not an algebraic ornament. It is the finite ledger's witness that order has
measurement content.

Why should this close only through successor counting? Because the grammar
cannot read an ordered difference without counting the order. The first
successor carries the first admissible transport. The second successor
carries the next transport. A further comparison reads whether the two-step
path can be identified with the path taken in reverse. Without the successor
discipline, the reader would speak as if \(AB\) and \(BA\) already lived in a
completed arena where subtraction could be applied freely. The grammar has
not granted that arena.

Successor counting is austere. It does not bring in all numbers at once. It
permits a next step, then another next step, and requires each step to say
what it carries. The commutator residue needs exactly that discipline. One
step alone cannot expose noncommutation. Two steps can create the ordered
comparison. The return comparison asks whether the two ordered readings close
on the same extensional condition. If they do not, the residue is the
unclosed debt.

The count also prevents a false cancellation. If the ledger forgot which
transport occurred first, \(AB\) and \(BA\) could be collapsed by presentation
instead of by measurement. Extensionality would then be misused. The
extensional quotient identifies conditions with the same carried obligations;
it does not erase order when order changes what is carried. The successor
count keeps the order recoverable, and therefore keeps the residue honest.

Closure means that the residue has been counted, bounded, and assigned to the
ledger. It does not mean the residue disappears by decree. A gauge
commutator closes when the grammar has enough successor structure to read
the two ordered paths, compare them extensionally, and carry any remaining
difference as a bounded residue. The close-out is finite even when it points
toward completion. It says: this is the residual debt produced by these
ordered transports, and this debt is now in the count.

A loop illustration can help. In a Mach--Zehnder reading, two admissible
branches may remain separately consistent until they meet at a shared
boundary. If their carried distinctions agree, the boundary fold identifies
them. If they return opposed, the fold cannot form a single ledger without
cancelling the inadmissible ordering. The example is optical in its familiar
setting, but the grammatical lesson is simpler: ordered branches can carry a
residue that is invisible to each branch alone and readable only at the
comparison.

The Aharonov--Bohm illustration says the same thing in gauge language. Local
segments may look silent with respect to the field reading, while the closed
transport around a loop returns a holonomy. The chapter does not need the
physical apparatus to establish the point. The point is grammatical: a
connection can carry memory around a comparison, and the commutator residue
is the finite mark that order has not quotiented away.

This is also why the count reaching three matters. The first count permits a
transport to be carried. The second count permits the ordered pair of
transports. The third count permits the comparison to return as a closed
ledger entry rather than as an untracked alternation. The grammar is not
celebrating the numeral. It is recording the minimum successor discipline
needed for a loop residue to become a stable finite obligation.

Once the commutator residue has closed in this way, it can be handed to
completion. But completion must not wash out its finite origin. The residue
was not born in a smooth space. It was born from extensional conditions and
successor-ordered transport. The next section therefore asks how this finite
residue maps into the Sobolev--Hilbert completion while preserving the split
that made the gauge reading lawful.
